\justifying
% === INTRODUCTION GÉNÉRALE ===
\section*{Introduction générale}
\addcontentsline{toc}{section}{Introduction générale}

Dans un environnement technologique en constante évolution, la maintenance des systèmes informatiques représente un enjeu stratégique majeur pour les entreprises de services numériques.

Les organisations doivent concilier la réactivité face aux incidents, le respect des accords de niveau de service SLA (\textit{Service Level Agreement}) et la traçabilité des interventions.
\newline

\textbf{SolidWall Consulting}, acteur reconnu dans le domaine du développement web et mobile, n'échappe pas à cette réalité. En tant que prestataire de services numériques gérant un portefeuille croissant de clients et de contrats de maintenance, l'entreprise se trouve confrontée à des exigences de plus en plus élevées en matière de réactivité, de traçabilité et de respect des engagements de service. La maîtrise de ces enjeux constitue pour elle un facteur clé de compétitivité et de fidélisation de sa clientèle.
\newline

C’est dans ce contexte que s’inscrit notre projet de fin d’études, visant à concevoir et réaliser \textbf{SolidMaint}, une plateforme web intégrée de gestion des contrats de maintenance. L’objectif principal consiste à automatiser et centraliser l’ensemble du processus métier, depuis la soumission des demandes clients jusqu’à la résolution des incidents, tout en garantissant une traçabilité complète et le respect des engagements contractuels.
\newline

Ce présent rapport s'articule autour de six chapitres retraçant la démarche complète du projet :

Le \textbf{premier chapitre} établit le cadre général avec l'analyse de l'existant, l'étude comparative des solutions du marché, et la définition de la solution proposée.

Le \textbf{deuxième chapitre} détaille la spécification des besoins fonctionnels et non fonctionnels, l'identification des acteurs du système, la modélisation UML des cas d'utilisation, la planification Scrum avec constitution du Product Backlog, l'environnement de développement, et l'architecture technique de l'application.

Le \textbf{troisième chapitre} présente la réalisation du premier sprint : authentification sécurisée, système rôles et permissions, gestion des utilisateurs et des clients.

Les \textbf{chapitres suivants} couvrent les sprints restants : gestion des contrats, processus des demandes et tickets, fonctionnalités de collaboration et intelligence artificielle, accompagnés des tests et validations.



